\documentclass{article}
\usepackage[preprint]{neurips_2025}
\usepackage[T1]{fontenc}
\usepackage{amsmath,amssymb,booktabs,tabularx}
\usepackage{microtype}
\usepackage[hidelinks]{hyperref}
\usepackage{url}
\hypersetup{pdftitle={Asymmetric Process-Reward Adaptation for Mathematical Reasoning},
  pdfauthor={Hongyi Ding; Jasmine Jiang; Yang Zeng}, pdfsubject={ECE285 Project Proposal}}
\title{Asymmetric Process-Reward Adaptation\\for Mathematical Reasoning}
\author{
Hongyi Ding \\
\And
Jasmine Jiang \\
\And
Yang Zeng \\
}
\newcommand{\E}{\mathbb{E}}
\newcommand{\sg}{\operatorname{sg}}
\newcolumntype{L}{>{\raggedright\arraybackslash}X}
\begin{document}
\maketitle
\begin{abstract}
Correct final answers can conceal faulty derivations, making outcome labels noisy supervision for process reward models (PRMs). We propose asymmetric offline adaptation: retain failure supervision and public step labels while omitting the positive outcome term. Four matched GRPO experiments test whether this change improves small-model mathematical reasoning.
\end{abstract}

\section{Motivation and research question}
Mathematical RL needs rewards that distinguish valid reasoning from lucky answers. Yet answer-correct solutions can contain intermediate errors \citep{zheng2025}, so positive outcome supervision may teach a PRM to certify flawed derivations. We ask: \textbf{can withholding this supervision improve downstream policy accuracy?} We hypothesize that less false certification can outweigh lost useful supervision, yielding better rewards without new annotations.

\section{Related work and contribution}
PRM800K supplies step labels \citep{lightman2023}; PRIME learns implicit rewards online \citep{cui2025}; PURE addresses reward hacking through min-form credit assignment \citep{cheng2025}. LCA learns PRMs with softmax-weighted pooling and already studies one-sided label noise \citep{jia2026}. Our symmetric control adopts its pooling. The proposed contribution is an \textbf{asymmetric supervision rule with retained step labels}, evaluated through downstream GRPO. The D--C comparison isolates the removed loss term; C is LCA-inspired, not a full reproduction.

\section{Proposed method}
\textbf{Offline adaptation.} Generate four solutions per problem on 512 training problems using the initial policy. Retain complete, unambiguously verified trajectories. For native step-boundary correctness scores $p_{\phi,t}$, set $e_{\phi,t}=1-p_{\phi,t}$. Following LCA,
\begin{equation}
 w_t=\frac{\exp(e_{\phi,t}/T)}{\sum_j\exp(e_{\phi,j}/T)},\qquad
 E_\phi(\tau)=\sum_t\sg(w_t)e_{\phi,t},\qquad T=0.2 .
\end{equation}
Here $\sg$ stops gradients through weights. Define $L_{\rm fail}=\E_{\tau^-}[-\log E_\phi(\tau^-)]$ and $L_{\rm success}=\E_{\tau^+}[-\log(1-E_\phi(\tau^+))]$, taking separate means over failures/successes. A shared set of 2,048 public positive/negative step labels supplies mean binary cross-entropy $L_{\rm step}$:
\begin{equation}
 L_c=L_{\rm fail}+cL_{\rm success}+L_{\rm step},\qquad
 c_C=1,\quad \boxed{c_D=0\ \text{(ours)}} .
\end{equation}
Failure labels provide weak evidence of an error, without locating it. Real step labels discourage blanket rejection; downstream gains remain an empirical question.

\newpage
\noindent\textbf{Policy training.} Freeze each PRM throughout RL. For complete, parseable responses in B--D, compute
\begin{equation}
 R_i=y_i+0.2\min_t p_{\phi,t}(\tau_i),\qquad
 \hat A_i=\frac{R_i-\bar R}{\sigma_R+10^{-4}} .
\end{equation}
Here $y_i\in\{0,1\}$ is rule-verified answer correctness; $\bar R,\sigma_R$ are the group mean and sample standard deviation. A uses $R_i=y_i$; truncated or unparseable responses receive zero reward in all arms. GRPO applies the detached sequence advantage to completion tokens \citep{shao2024}. \textbf{Correct policy rollouts remain in RL and retain their outcome reward.}

\section{Experimental design and expected findings}
\begin{center}
\small
\renewcommand{\arraystretch}{1.12}
\begin{tabularx}{\linewidth}{@{}l L L@{}}
\toprule
\textbf{Arm} & \textbf{Reward and PRM preparation} & \textbf{Question tested} \\
\midrule
A & Outcome only; no PRM & Reference GRPO performance \\
B & Outcome + original frozen PRM & B--A: value of an existing PRM \\
C & Outcome + PRM adapted with $c=1$ & C--B: value of full adaptation \\
\textbf{D (ours)} & Outcome + PRM adapted with $c=0$ & \textbf{D--C: value of the proposed change} \\
\bottomrule
\end{tabularx}
\end{center}
\textbf{Models and implementation.} Use Qwen2.5-Math-1.5B as policy and the H4 1.5B PRM \citep{h4prm}, with LoRA ranks 16 and 8. TRL 0.23.1/PEFT uses \texttt{loss\_type="grpo"} and sums reward components before group normalization \citep{trl}. The outcome-only training and evaluation pipeline has been verified on one 80\,GB A100; PRM-based comparisons remain to be run.

\textbf{Data and contamination.} The audited MATH split contains 6,993 training and 500 validation problems \citep{hendrycks2021}. Use \textbf{MATH500} for primary evaluation and GSM8K test for transfer \citep{cobbe2021}. Neither has exact normalized question overlap with the public PRM800K training inventory. All newly selected anchors will exclude validation, test, and diagnostic questions. Since PRM800K incorporates 4,500 original MATH test questions \citep{prmdata}, full MATH is supplementary only. Unknown pretraining exposure and paraphrases remain limitations.

\textbf{Matched controls and budget.} C/D share PRM initialization, batches, optimizer, and update count; only $c$ changes, without re-averaging the remaining losses. All arms share policy initialization, prompt order, decoding, and tuning effort: eight rollouts/question, temperature 1.0, and a 2,048-token cap. Run four arms with three paired seeds at \textbf{24M generated completion tokens/run} (288M total), retaining 6M/12M/24M milestones. Count incorrect and truncated completions; report update counts, prompt counts, response lengths, and end-to-end GPU hours, including PRM costs.

\textbf{Evaluation and interpretation.} Evaluate the initial policy and all three milestones with identical greedy decoding, without PRM reranking or voting. The primary comparison is \textbf{D--C MATH500 accuracy at 24M tokens}; D--A/B measures practical value. Report per-seed differences, mean/std, and 95\% paired question-cluster bootstrap intervals, keeping all arm/seed predictions for a question together. Intermediate scores measure learning progress and do not select the final checkpoint.

\textbf{Mechanism and risks.} On fixed ProcessBench MATH solutions, measure first-error accuracy, correct-chain acceptance, and their harmonic mean at threshold 0.5. Within answer-correct solutions, valid-chain AUROC tests whether the PRM distinguishes valid from flawed reasoning beyond a uniform score shift. Report known overlaps and class counts; use no diagnostic labels for tuning. Benefits require downstream gains: lower PRM scores alone are insufficient, and small or inconsistent accuracy changes remain inconclusive. Reward hacking and policy distribution shift are residual risks.

\textbf{Tentative timeline.} Week 1: finalize the offline buffer, anchor audit, and PRM adaptation. Week 2: complete A--D for the first seed. Week 3: run the remaining paired seeds and reward diagnostics. Week 4: complete statistical analysis, error analysis, and the reproducible report.
\clearpage
\begin{thebibliography}{11}
\bibitem[Shao et al.(2024)]{shao2024}
Zhihong Shao et al. DeepSeekMath: Pushing the Limits of Mathematical Reasoning in Open Language Models. 2024.
\url{https://arxiv.org/abs/2402.03300v3}.

\bibitem[Zheng et al.(2025)]{zheng2025}
Chujie Zheng et al. ProcessBench: Identifying Process Errors in Mathematical Reasoning. ACL, 2025.
\url{https://arxiv.org/abs/2412.06559v4}.

\bibitem[Lightman et al.(2023)]{lightman2023}
Hunter Lightman et al. Let's Verify Step by Step. 2023.
\url{https://arxiv.org/abs/2305.20050}.

\bibitem[Cui et al.(2025)]{cui2025}
Ganqu Cui et al. Process Reinforcement through Implicit Rewards. 2025.
\url{https://arxiv.org/abs/2502.01456v2}.

\bibitem[Cheng et al.(2025)]{cheng2025}
Jie Cheng et al. Stop Summation: Min-Form Credit Assignment Is All Process Reward Model Needs for Reasoning. NeurIPS, 2025.
\url{https://arxiv.org/abs/2504.15275v3}.

\bibitem[Jia et al.(2026)]{jia2026}
Tianyu Jia et al. The Weakest Link Tells It All: Outcome-Supervised Process Reward Modeling via Learnable Credit Assignment. EMNLP, 2026. Camera-ready version.
\url{https://arxiv.org/abs/2606.27739v2}.

\bibitem[Hendrycks et al.(2021)]{hendrycks2021}
Dan Hendrycks et al. Measuring Mathematical Problem Solving With the MATH Dataset. NeurIPS, 2021.
\url{https://arxiv.org/abs/2103.03874v2}.

\bibitem[Cobbe et al.(2021)]{cobbe2021}
Karl Cobbe et al. Training Verifiers to Solve Math Word Problems. 2021.
\url{https://arxiv.org/abs/2110.14168v2}.

\bibitem[Hugging Face H4(2025)]{h4prm}
Hugging Face H4. Qwen2.5-Math-1.5B-Instruct-PRM-0.2. 2025.
\href{https://huggingface.co/HuggingFaceH4/Qwen2.5-Math-1.5B-Instruct-PRM-0.2}{Official model card and scoring interface}. Accessed Oct.\ 20, 2026.

\bibitem[Hugging Face(2026)]{trl}
Hugging Face. TRL 0.23.1: \href{https://huggingface.co/docs/trl/v0.23.1/grpo_trainer}{GRPO Trainer documentation}. Accessed Oct.\ 20, 2026.

\bibitem[OpenAI(2023)]{prmdata}
OpenAI. PRM800K dataset repository: MATH splits and answer grading. 2023. Accessed Oct.\ 20, 2026.
\url{https://github.com/openai/prm800k}.
\end{thebibliography}
\end{document}
